% Figure: the two normal modes of a symmetric two-mass three-spring chain.
% Top: the physical system, two equal masses joined by a coupling spring and
% anchored by two equal outer springs. Below: the in-phase mode (both masses
% move together, the coupling spring never stretches, low frequency) and the
% out-of-phase mode (masses move oppositely, the coupling spring works hardest,
% high frequency). Arrows show the instantaneous displacement of each mass.
\begin{figure}[htbp]
\centering
\begin{tikzpicture}[
  wall/.style={draw=engblack, thick},
  spring/.style={draw=enggray, decorate,
    decoration={coil, aspect=0.5, segment length=3pt, amplitude=3pt}},
  mass/.style={draw=engblack, fill=engblue!18, thick, minimum width=12mm,
    minimum height=8mm},
  arrowr/.style={draw=engred, -{Stealth}, very thick},
  label/.style={font=\small},
]
% ===== physical system (top) =====
\node[font=\small, anchor=west] at (0,3.4) {physical system:};
\draw[wall] (0.2,2.2) -- (0.2,3.1);
\draw[wall] (8.8,2.2) -- (8.8,3.1);
\draw[spring] (0.2,2.65) -- (2.1,2.65);
\node[mass] (m1) at (2.7,2.65) {$m$};
\draw[spring] (3.3,2.65) -- (5.5,2.65);
\node[mass] (m2) at (6.1,2.65) {$m$};
\draw[spring] (6.7,2.65) -- (8.8,2.65);
\node[font=\scriptsize, enggray] at (1.15,3.0) {$k$};
\node[font=\scriptsize, enggray] at (4.4,3.0) {$k_c$};
\node[font=\scriptsize, enggray] at (7.75,3.0) {$k$};

% ===== in-phase mode (middle) =====
\node[font=\small, anchor=west] at (0,1.55)
  {in-phase mode $\;\omega_1=\sqrt{k/m}$:};
\node[mass, minimum width=10mm, minimum height=6mm] (a1) at (2.7,0.8) {};
\node[mass, minimum width=10mm, minimum height=6mm] (a2) at (6.1,0.8) {};
\draw[arrowr] (a1.east) -- ++(0.9,0);
\draw[arrowr] (a2.east) -- ++(0.9,0);
\node[font=\scriptsize, enggray, anchor=west] at (7.2,0.8)
  {coupling spring $k_c$ never stretches};

% ===== out-of-phase mode (bottom) =====
\node[font=\small, anchor=west] at (0,-0.35)
  {out-of-phase mode $\;\omega_2=\sqrt{(k+2k_c)/m}$:};
\node[mass, minimum width=10mm, minimum height=6mm] (b1) at (2.7,-1.1) {};
\node[mass, minimum width=10mm, minimum height=6mm] (b2) at (6.1,-1.1) {};
\draw[arrowr] (b1.east) -- ++(0.9,0);
\draw[arrowr] (b2.west) -- ++(-0.9,0);
\node[font=\scriptsize, enggray, anchor=west] at (7.2,-1.1)
  {coupling spring $k_c$ works hardest};
\end{tikzpicture}
\caption[Normal modes of a two-mass coupled oscillator]{A symmetric chain of
two equal masses, two outer anchoring springs of stiffness $k$, and a central
coupling spring $k_c$ has two normal modes read straight off its energy. In
the in-phase mode both masses move together, the coupling spring keeps its
natural length and stores no energy, and the frequency is that of a single
outer spring alone, $\omega_1=\sqrt{k/m}$. In the out-of-phase mode the masses
move oppositely, the coupling spring is stretched to twice the mass
displacement and contributes its full stiffness, raising the frequency to
$\omega_2=\sqrt{(k+2k_c)/m}$. A general small motion is a superposition of the
two, and the slow beat between them is the energy trading back and forth
between the masses.}
\label{fig:vol03ch05:two-dof-modes}
\end{figure}
